\section{Background and Motivation}
\label{sec:motivation}
\subsection{Recovery work and workflow delay differ}
A workflow is a DAG of requests, tools and control nodes. A join waits for
every predecessor, so a recovery pause on one branch can be absorbed while
another branch runs. Figure~\ref{fig:motivation} connects observed branch
imbalance to this mechanism using the public CC/Weka trace~\cite{semianalysis2026cctraces}.
We retain parser fan-outs whose branch intervals have nonempty common
overlap, and take one maximum longest-minus-shortest duration gap per
workflow. The 124 qualifying workflows have median and 90th-percentile gaps
of 140.6 and 1,570.3 s. These are duration differences, not duration ratios or
unbiased estimates of every branch's available Slack.

\begin{figure*}[tp]
\centering\includegraphics[width=\textwidth]{figures/motivation-real-cc.pdf}
\caption{Observed imbalance and modeled recovery consequence.
(a) One maximum strict-parallel duration gap per qualifying CC/Weka workflow.
(b) The same five-branch trace example; solid bars are observed execution,
dashed segments are time remaining until the reconstructed join. Parent edges
are inferred from ordering and timestamps.
(c) A separate fault on each branch at 349.5 s: modeled recomputation is split
into absorbed and exposed delay; labels give exposed seconds. This frozen
AIConfigurator counterfactual uses Qwen3-235B-A22B, A100-80GB, TP8/PP4/EP8;
it is neither a measured failure nor the Qwen3-8B evaluation.}
\Description{A workflow-level empirical CDF, a time-scaled five-branch
execution diagram with join slack, and a stacked comparison of absorbed
versus exposed modeled recovery delay.}
\label{fig:motivation}
\end{figure*}

The example separates runtime length, join slack and recovery cost. A3 has
52.4 s before the join and absorbs its modeled 46.6 s recomputation. A4 is
critical and exposes all 35.9 s; A2 and A5 expose intermediate delays.
The frozen counterfactual~\cite{nvidia2026aiconfigurator} substitutes a configured
recomputation model for the trace's original model. Input lengths come from
64-token hash-block counts, not exact retokenization. It assumes unchanged
other branches and omits recovery contention; it illustrates a mechanism
without estimating deployment speedup. The later simulation uses a separate,
exact-token FanOutQA-derived workload.

The local harm approximation $J_i=\pos{T_i-s_i}$ relates recovery time $T_i$
to remaining Slack $s_i$. Queues, contention and simultaneous failures can
change the critical branch. AegisKV therefore refreshes Slack during execution,
and evaluation measures complete workflows.

\subsection{KV protection consumes the resources it aims to save}
For a model with $L$ attention layers, $h_{kv}$ KV heads, head dimension $d$
and $b$ bytes per element, a token occupies $2Lh_{kv}db$ KV bytes before
allocator rounding. Our Qwen3-8B layout uses 147,456 bytes per token: a
4,096-token prompt chunk occupies 576 MiB, and an 86,016-token prompt roughly
11.81 GiB. Figure~\ref{fig:motivation-cost} quantifies the footprint and separates it
from measured pilot transfer rates. Two copies of a long prompt can exceed the memory left
after weights and runtime reserves on a modeled 40 GiB GPU.

HBM protection competes with primary KV; DRAM requires source-to-host and
host-to-recovery copies. Both consume serving resources. The pilot restores
HBM state yet runs slightly slower (Section~\ref{sec:physical}).

\subsection{Requirements}
Protection must value workflow harm avoided by recoverable prefixes. Each
tier must remain independently useful, regardless of the other's residency.
Physical state must govern replacement and recovery: pending copies are
unusable, and committed backups retain storage while populating primary KV.

\begin{figure*}[tp]
\centering\includegraphics[width=\textwidth]{figures/aligned-motivation-cost.pdf}
\caption{Capacity and service cost. (a) Analytical Qwen3-8B KV footprint
and modeled post-weight/runtime HBM budget; these are not measured allocations.
(b) Every physical RTX4090 calibration observation (circles) and its
aggregate byte/time rate (diamonds); $n$ is the number of transfer samples.
Host staging and software overhead are included. One recovery sample per
path supports neither an uncertainty interval nor A100 calibration.}
\Description{One- and two-copy KV footprints versus prompt length, alongside
measured HBM and DRAM protection/recovery service rates in the physical pilot.}
\label{fig:motivation-cost}
\end{figure*}
